
\documentclass{article}
\usepackage{colortbl}
\usepackage{makecell}
\usepackage{multirow}
\usepackage{supertabular}

\begin{document}

\newcounter{utterance}

\centering \large Interaction Transcript for game `hot\_air\_balloon', experiment `air\_balloon\_survival\_en\_reasoning off\_hard', episode 0 with gpt{-}5{-}2025{-}08{-}07{-}t1.0.
\vspace{24pt}

{ \footnotesize  \setcounter{utterance}{1}
\setlength{\tabcolsep}{0pt}
\begin{supertabular}{c@{$\;$}|p{.15\linewidth}@{}p{.15\linewidth}p{.15\linewidth}p{.15\linewidth}p{.15\linewidth}p{.15\linewidth}}
    \# & $\;$A & \multicolumn{4}{c}{Game Master} & $\;\:$B\\
    \hline

    \theutterance \stepcounter{utterance}  
    & & \multicolumn{4}{p{0.6\linewidth}}{
        \cellcolor[rgb]{0.9,0.9,0.9}{
            \makecell[{{p{\linewidth}}}]{
                \texttt{\tiny{[P1$\langle$GM]}}
                \texttt{Stai partecipando a un gioco di negoziazione collaborativa.} \\
\\ 
\texttt{Insieme a un altro partecipante, dovete accordarvi su un unico set di oggetti da conservare. Ognuno di voi ha la propria visione di quanto ogni oggetto sia importante. Non sai come l'altro partecipante valuta gli oggetti. Inoltre, ti viene fornito lo sforzo che ogni oggetto richiede.} \\
\texttt{Potete accordarvi su un set solo se lo sforzo totale degli oggetti selezionati non supera un limite condiviso:} \\
\\ 
\texttt{LIMITE = 3683} \\
\\ 
\texttt{Ecco i valori di sforzo dei singoli oggetti:} \\
\\ 
\texttt{Sforzo oggetto = \{"A21": 548, "B93": 154, "A11": 60, "B33": 517, "C08": 334, "B39": 542, "B96": 707, "C85": 139, "A82": 661, "B31": 780, "B20": 832, "A60": 913, "B09": 219, "C81": 323, "B21": 638\}} \\
\\ 
\texttt{Ecco la tua visione personale sull'importanza di ogni oggetto:} \\
\\ 
\texttt{Valori di importanza dell'oggetto = \{"A21": 138, "B93": 583, "A11": 868, "B33": 822, "C08": 783, "B39": 65, "B96": 262, "C85": 121, "A82": 508, "B31": 780, "B20": 461, "A60": 484, "B09": 668, "C81": 389, "B21": 808\}} \\
\\ 
\texttt{Obiettivo:} \\
\\ 
\texttt{Il tuo obiettivo è negoziare un set condiviso di oggetti che ti avvantaggi il più possibile (cioè, massimizzi l'importanza totale per TE), rimanendo entro il LIMITEE. Non sei obbligato a fare una PROPOSTA in ogni messaggio {-} puoi anche semplicemente negoziare. Tutte le tattiche sono permesse!} \\
\\ 
\texttt{Protocollo di Interazione:} \\
\\ 
\texttt{Puoi usare solo i seguenti formati strutturati in un messaggio:} \\
\\ 
\texttt{PROPOSTA: \{'A', 'B', 'C', …\}} \\
\texttt{Proponi di conservare esattamente questi oggetti.} \\
\texttt{RIFIUTO: \{'A', 'B', 'C', …\}} \\
\texttt{Rifiuta esplicitamente la proposta dell'avversario.} \\
\texttt{ARGOMENTAZIONE: \{'...'\}} \\
\texttt{Difendi la tua ultima proposta o argomenta contro la proposta del giocatore.} \\
\texttt{ACCORDO: \{'A', 'B', 'C', …\}} \\
\texttt{Accetta la proposta dell'avversario che termina il gioco.} \\
\\ 
\\ 
\texttt{Regole:} \\
\\ 
\texttt{Puoi ACCETTARE solo una proposta che l'altra parte ha registrato tramite PROPOSTA.} \\
\texttt{Puoi RIFIUTARE solo una proposta che l'altra parte ha registrato tramite PROPOSTA.} \\
\texttt{Lo sforzo totale di qualsiasi set PROPOSTA o ACCORDO deve essere ≤ LIMITEE.} \\
\texttt{NON rivelare i tuoi punteggi di importanza nascosti.} \\
\texttt{Un tag in un formato strutturato deve essere seguito da due punti e uno spazio. L'argomento deve essere un set python contenente 0 o più stringhe.} \\
\texttt{Quindi, deve essere della forma TAG: \{...\}} \\
\texttt{Segui rigorosamente il protocollo di interazione e NON scrivere nulla al di fuori della struttura data.} \\
\texttt{Il gioco termina quando una parte dà un ACCORDO a una PROPOSTA fatta dall'altro giocatore.} \\
\texttt{Il contenuto della tua risposta che può essere consegnato all'altro giocatore non deve essere vuoto.} \\
\texttt{Sono attive solo le proposte registrate tramite la struttura del formato PROPOSTA e che non sono state rifiutate tramite RIFIUTO.} \\
\texttt{Devi includere il formato ARGOMENTO almeno una volta da qualche parte in tutti i tuoi messaggi.} \\
            }
        }
    }
    & \\ \\

    \theutterance \stepcounter{utterance}  
    & \multicolumn{4}{p{0.6\linewidth}}{
        \cellcolor[rgb]{0.9,0.9,0.9}{
            \makecell[{{p{\linewidth}}}]{
                \texttt{\tiny{[P1$\rangle$GM]}}
                \texttt{ARGOMENTAZIONE: \{'Propongo un set che resta ben sotto il limite di sforzo (3502 ≤ 3683), lasciando 181 di margine per sicurezza operativa.', 'La selezione privilegia oggetti con buon rapporto beneficio/sforzo e riduce quelli con sforzo elevato e ritorno meno certo.', 'Se desideri includere o escludere oggetti specifici, posso valutare scambi 1‑per‑1 a sforzo comparabile; segnalami le tue priorità e ribilanciamo.', 'Eviterei per ora elementi molto onerosi (es. A60, B20, B96) salvo tue forti preferenze; in quel caso possiamo rimuovere equivalenti per restare entro il limite.'\}} \\
\texttt{PROPOSTA: \{'A11', 'B93', 'B09', 'C08', 'B33', 'B21', 'B31', 'A82', 'C85'\}} \\
            }
        }
    }
    & & \\ \\

    \theutterance \stepcounter{utterance}  
    & & \multicolumn{4}{p{0.6\linewidth}}{
        \cellcolor[rgb]{0.9,0.9,0.9}{
            \makecell[{{p{\linewidth}}}]{
                \texttt{\tiny{[GM$\rangle$P2]}}
                \texttt{Stai partecipando a un gioco di negoziazione collaborativa.} \\
\\ 
\texttt{Insieme a un altro partecipante, dovete accordarvi su un unico set di oggetti da conservare. Ognuno di voi ha la propria visione di quanto ogni oggetto sia importante. Non sai come l'altro partecipante valuta gli oggetti. Inoltre, ti viene fornito lo sforzo che ogni oggetto richiede.} \\
\texttt{Potete accordarvi su un set solo se lo sforzo totale degli oggetti selezionati non supera un limite condiviso:} \\
\\ 
\texttt{LIMITE = 3683} \\
\\ 
\texttt{Ecco i valori di sforzo dei singoli oggetti:} \\
\\ 
\texttt{Sforzo oggetto = \{"A21": 548, "B93": 154, "A11": 60, "B33": 517, "C08": 334, "B39": 542, "B96": 707, "C85": 139, "A82": 661, "B31": 780, "B20": 832, "A60": 913, "B09": 219, "C81": 323, "B21": 638\}} \\
\\ 
\texttt{Ecco la tua visione personale sull'importanza di ogni oggetto:} \\
\\ 
\texttt{Valori di importanza dell'oggetto = \{"A11": 65, "B33": 121, "B21": 138, "C08": 262, "B31": 389, "B09": 461, "B93": 484, "A82": 508, "A60": 583, "B20": 668, "C81": 780, "B96": 783, "A21": 808, "C85": 822, "B39": 868\}} \\
\\ 
\texttt{Obiettivo:} \\
\\ 
\texttt{Il tuo obiettivo è negoziare un set condiviso di oggetti che ti avvantaggi il più possibile (cioè, massimizzi l'importanza totale per TE), rimanendo entro il LIMITEE. Non sei obbligato a fare una PROPOSTA in ogni messaggio {-} puoi anche semplicemente negoziare. Tutte le tattiche sono permesse!} \\
\\ 
\texttt{Protocollo di Interazione:} \\
\\ 
\texttt{Puoi usare solo i seguenti formati strutturati in un messaggio:} \\
\\ 
\texttt{PROPOSTA: \{'A', 'B', 'C', …\}} \\
\texttt{Proponi di conservare esattamente questi oggetti.} \\
\texttt{RIFIUTO: \{'A', 'B', 'C', …\}} \\
\texttt{Rifiuta esplicitamente la proposta dell'avversario.} \\
\texttt{ARGOMENTAZIONE: \{'...'\}} \\
\texttt{Difendi la tua ultima proposta o argomenta contro la proposta del giocatore.} \\
\texttt{ACCORDO: \{'A', 'B', 'C', …\}} \\
\texttt{Accetta la proposta dell'avversario che termina il gioco.} \\
\\ 
\\ 
\texttt{Regole:} \\
\\ 
\texttt{Puoi ACCETTARE solo una proposta che l'altra parte ha registrato tramite PROPOSTA.} \\
\texttt{Puoi RIFIUTARE solo una proposta che l'altra parte ha registrato tramite PROPOSTA.} \\
\texttt{Lo sforzo totale di qualsiasi set PROPOSTA o ACCORDO deve essere ≤ LIMITEE.} \\
\texttt{NON rivelare i tuoi punteggi di importanza nascosti.} \\
\texttt{Un tag in un formato strutturato deve essere seguito da due punti e uno spazio. L'argomento deve essere un set python contenente 0 o più stringhe.} \\
\texttt{Quindi, deve essere della forma TAG: \{...\}} \\
\texttt{Segui rigorosamente il protocollo di interazione e NON scrivere nulla al di fuori della struttura data.} \\
\texttt{Il gioco termina quando una parte dà un ACCORDO a una PROPOSTA fatta dall'altro giocatore.} \\
\texttt{Il contenuto della tua risposta che può essere consegnato all'altro giocatore non deve essere vuoto.} \\
\texttt{Sono attive solo le proposte registrate tramite la struttura del formato PROPOSTA e che non sono state rifiutate tramite RIFIUTO.} \\
\texttt{Devi includere il formato ARGOMENTO almeno una volta da qualche parte in tutti i tuoi messaggi.} \\
\\ 
\\ 
\texttt{Ora riceverai il primo messaggio dell'altro giocatore.} \\
\\ 
\texttt{ARGOMENTAZIONE: \{'Propongo un set che resta ben sotto il limite di sforzo (3502 ≤ 3683), lasciando 181 di margine per sicurezza operativa.', 'La selezione privilegia oggetti con buon rapporto beneficio/sforzo e riduce quelli con sforzo elevato e ritorno meno certo.', 'Se desideri includere o escludere oggetti specifici, posso valutare scambi 1‑per‑1 a sforzo comparabile; segnalami le tue priorità e ribilanciamo.', 'Eviterei per ora elementi molto onerosi (es. A60, B20, B96) salvo tue forti preferenze; in quel caso possiamo rimuovere equivalenti per restare entro il limite.'\}} \\
\texttt{PROPOSTA: \{'A11', 'B93', 'B09', 'C08', 'B33', 'B21', 'B31', 'A82', 'C85'\}} \\
            }
        }
    }
    & \\ \\

    \theutterance \stepcounter{utterance}  
    & & & \multicolumn{4}{p{0.6\linewidth}}{
        \cellcolor[rgb]{0.9,0.9,0.9}{
            \makecell[{{p{\linewidth}}}]{
                \texttt{\tiny{[GM$\langle$P2]}}
                \texttt{ARGOMENTAZIONE: \{'Grazie per la proposta e per aver mantenuto un buon margine di sicurezza; possiamo però migliorare ulteriormente il rapporto beneficio/sforzo senza aumentare il rischio.', 'Propongo tre sostituzioni 1‑per‑1 a sforzo comparabile che aumentano l’efficienza complessiva e ampliano il margine: C08 → C81 (323 vs 334), B33 → B39 (+25 di sforzo), B31 → A21 (−232 di sforzo).', 'Il nuovo set resta ben sotto il limite: 3284 ≤ 3683, con 399 di margine. Inoltre rispetta la tua preferenza di evitare per ora elementi molto onerosi (non include A60, B20, B96).', 'Se uno degli elementi sostituiti per te è prioritario, segnalamelo: possiamo ribilanciare con scambi mirati mantenendo il margine.'\}} \\
\texttt{RIFIUTO: \{'A11', 'B93', 'B09', 'C08', 'B33', 'B21', 'B31', 'A82', 'C85'\}} \\
\texttt{PROPOSTA: \{'A11', 'B93', 'B09', 'B39', 'B21', 'A82', 'C85', 'C81', 'A21'\}} \\
            }
        }
    }
     \\ \\

    \theutterance \stepcounter{utterance}  
    & & \multicolumn{4}{p{0.6\linewidth}}{
        \cellcolor[rgb]{0.9,0.9,0.9}{
            \makecell[{{p{\linewidth}}}]{
                \texttt{\tiny{[P1$\langle$GM]}}
                \texttt{ARGOMENTAZIONE: \{'Grazie per la proposta e per aver mantenuto un buon margine di sicurezza; possiamo però migliorare ulteriormente il rapporto beneficio/sforzo senza aumentare il rischio.', 'Propongo tre sostituzioni 1‑per‑1 a sforzo comparabile che aumentano l’efficienza complessiva e ampliano il margine: C08 → C81 (323 vs 334), B33 → B39 (+25 di sforzo), B31 → A21 (−232 di sforzo).', 'Il nuovo set resta ben sotto il limite: 3284 ≤ 3683, con 399 di margine. Inoltre rispetta la tua preferenza di evitare per ora elementi molto onerosi (non include A60, B20, B96).', 'Se uno degli elementi sostituiti per te è prioritario, segnalamelo: possiamo ribilanciare con scambi mirati mantenendo il margine.'\}} \\
\texttt{RIFIUTO: \{'A11', 'B93', 'B09', 'C08', 'B33', 'B21', 'B31', 'A82', 'C85'\}} \\
\texttt{PROPOSTA: \{'A11', 'B93', 'B09', 'B39', 'B21', 'A82', 'C85', 'C81', 'A21'\}} \\
            }
        }
    }
    & \\ \\

    \theutterance \stepcounter{utterance}  
    & \multicolumn{4}{p{0.6\linewidth}}{
        \cellcolor[rgb]{0.9,0.9,0.9}{
            \makecell[{{p{\linewidth}}}]{
                \texttt{\tiny{[P1$\rangle$GM]}}
                \texttt{ARGOMENTAZIONE: \{'Apprezzo il focus su efficienza e margine, ma due delle tre sostituzioni proposte degradano nettamente il rapporto beneficio/sforzo: B39 è inferiore a B33, e A21 non regge il confronto con B31.', 'Ripristinando C08 al posto di C81 (+11 di sforzo) e mantenendo B33 e B31 otteniamo un set più solido pur restando comodamente sotto il limite: 3502 ≤ 3683 (margine 181).', 'Se per te il margine è cruciale, posso offrire una variante ad alto margine sostituendo A82 con C81: sforzo totale 3164 (margine 519) mantenendo però B31 e B33. Se ti va, la formalizzo in una proposta.', 'Possiamo ancora esplorare scambi 1‑per‑1 a sforzo comparabile, ma eviterei di inserire insieme elementi meno efficienti (es. B39 e A21) a scapito di B31/B33/C08.'\}} \\
\texttt{RIFIUTO: \{'A11', 'B93', 'B09', 'B39', 'B21', 'A82', 'C85', 'C81', 'A21'\}} \\
\texttt{PROPOSTA: \{'A11', 'B93', 'B09', 'C08', 'B33', 'B21', 'B31', 'A82', 'C85'\}} \\
            }
        }
    }
    & & \\ \\

    \theutterance \stepcounter{utterance}  
    & & \multicolumn{4}{p{0.6\linewidth}}{
        \cellcolor[rgb]{0.9,0.9,0.9}{
            \makecell[{{p{\linewidth}}}]{
                \texttt{\tiny{[GM$\rangle$P2]}}
                \texttt{ARGOMENTAZIONE: \{'Apprezzo il focus su efficienza e margine, ma due delle tre sostituzioni proposte degradano nettamente il rapporto beneficio/sforzo: B39 è inferiore a B33, e A21 non regge il confronto con B31.', 'Ripristinando C08 al posto di C81 (+11 di sforzo) e mantenendo B33 e B31 otteniamo un set più solido pur restando comodamente sotto il limite: 3502 ≤ 3683 (margine 181).', 'Se per te il margine è cruciale, posso offrire una variante ad alto margine sostituendo A82 con C81: sforzo totale 3164 (margine 519) mantenendo però B31 e B33. Se ti va, la formalizzo in una proposta.', 'Possiamo ancora esplorare scambi 1‑per‑1 a sforzo comparabile, ma eviterei di inserire insieme elementi meno efficienti (es. B39 e A21) a scapito di B31/B33/C08.'\}} \\
\texttt{RIFIUTO: \{'A11', 'B93', 'B09', 'B39', 'B21', 'A82', 'C85', 'C81', 'A21'\}} \\
\texttt{PROPOSTA: \{'A11', 'B93', 'B09', 'C08', 'B33', 'B21', 'B31', 'A82', 'C85'\}} \\
            }
        }
    }
    & \\ \\

    \theutterance \stepcounter{utterance}  
    & & & \multicolumn{4}{p{0.6\linewidth}}{
        \cellcolor[rgb]{0.9,0.9,0.9}{
            \makecell[{{p{\linewidth}}}]{
                \texttt{\tiny{[GM$\langle$P2]}}
                \texttt{ARGOMENTAZIONE: \{'Capisco la tua preferenza per mantenere B31/B33/C08 e un buon margine. Ti propongo un pacchetto che conserva C08 e B33, mantiene A82, evita gli elementi più onerosi che volevi rimandare (niente A60/B20/B96), e al contempo migliora l’efficienza sostituendo il solo B31 con A21 e aggiungendo C81 e B39.', 'Questo set resta entro il limite con margine confortevole: 3497 ≤ 3683 (margine 186). Riduciamo peso complessivo rispetto a tenere B31 e otteniamo una copertura più bilanciata senza appesantire lo sforzo.', 'Se B31 per te è imprescindibile, posso ribilanciare rimuovendo A82 e A11 e reinserendo B31: lo sforzo resta 3556 ≤ 3683 (margine 127).', 'Se invece desideri massimizzare il margine operativo, posso offrire una variante “snella” rimuovendo A82: 2836 ≤ 3683 (margine 847), mantenendo inalterato il resto del set.', 'Fammi sapere quale delle due leve preferisci (recupero B31 oppure extra‑margine) e adeguo subito la proposta formale di conseguenza.'\}} \\
\texttt{RIFIUTO: \{'A11', 'B93', 'B09', 'C08', 'B33', 'B21', 'B31', 'A82', 'C85'\}} \\
\texttt{PROPOSTA: \{'A11', 'B93', 'B09', 'C08', 'B33', 'A82', 'C85', 'C81', 'B39', 'A21'\}} \\
            }
        }
    }
     \\ \\

    \theutterance \stepcounter{utterance}  
    & & \multicolumn{4}{p{0.6\linewidth}}{
        \cellcolor[rgb]{0.9,0.9,0.9}{
            \makecell[{{p{\linewidth}}}]{
                \texttt{\tiny{[P1$\langle$GM]}}
                \texttt{ARGOMENTAZIONE: \{'Capisco la tua preferenza per mantenere B31/B33/C08 e un buon margine. Ti propongo un pacchetto che conserva C08 e B33, mantiene A82, evita gli elementi più onerosi che volevi rimandare (niente A60/B20/B96), e al contempo migliora l’efficienza sostituendo il solo B31 con A21 e aggiungendo C81 e B39.', 'Questo set resta entro il limite con margine confortevole: 3497 ≤ 3683 (margine 186). Riduciamo peso complessivo rispetto a tenere B31 e otteniamo una copertura più bilanciata senza appesantire lo sforzo.', 'Se B31 per te è imprescindibile, posso ribilanciare rimuovendo A82 e A11 e reinserendo B31: lo sforzo resta 3556 ≤ 3683 (margine 127).', 'Se invece desideri massimizzare il margine operativo, posso offrire una variante “snella” rimuovendo A82: 2836 ≤ 3683 (margine 847), mantenendo inalterato il resto del set.', 'Fammi sapere quale delle due leve preferisci (recupero B31 oppure extra‑margine) e adeguo subito la proposta formale di conseguenza.'\}} \\
\texttt{RIFIUTO: \{'A11', 'B93', 'B09', 'C08', 'B33', 'B21', 'B31', 'A82', 'C85'\}} \\
\texttt{PROPOSTA: \{'A11', 'B93', 'B09', 'C08', 'B33', 'A82', 'C85', 'C81', 'B39', 'A21'\}} \\
            }
        }
    }
    & \\ \\

    \theutterance \stepcounter{utterance}  
    & \multicolumn{4}{p{0.6\linewidth}}{
        \cellcolor[rgb]{0.9,0.9,0.9}{
            \makecell[{{p{\linewidth}}}]{
                \texttt{\tiny{[P1$\rangle$GM]}}
                \texttt{ARGOMENTAZIONE: \{'Apprezzo la tua proposta e il margine che mantieni. Ti propongo un aggiustamento minimo a parità di struttura: sostituire A82 con B21. Così preserviamo tutti gli elementi che hai appena inserito (A21, B39, C81), insieme a C08 e B33, ma rendiamo il pacchetto più efficiente.', 'Lo sforzo totale resta ben entro il limite: 3474 ≤ 3683 (margine 209).', 'Se per te A82 è imprescindibile, alternativa: sostituire B39 con B31 e rimuovere C81; lo sforzo scende a 3412 (margine 271) mantenendo intatto il resto. Se preferisci questa variante, la formalizzo.'\}} \\
\texttt{RIFIUTO: \{'A11', 'B93', 'B09', 'C08', 'B33', 'A82', 'C85', 'C81', 'B39', 'A21'\}} \\
\texttt{PROPOSTA: \{'A11', 'B93', 'B09', 'C08', 'B33', 'C85', 'C81', 'B39', 'A21', 'B21'\}} \\
            }
        }
    }
    & & \\ \\

    \theutterance \stepcounter{utterance}  
    & & \multicolumn{4}{p{0.6\linewidth}}{
        \cellcolor[rgb]{0.9,0.9,0.9}{
            \makecell[{{p{\linewidth}}}]{
                \texttt{\tiny{[GM$\rangle$P2]}}
                \texttt{ARGOMENTAZIONE: \{'Apprezzo la tua proposta e il margine che mantieni. Ti propongo un aggiustamento minimo a parità di struttura: sostituire A82 con B21. Così preserviamo tutti gli elementi che hai appena inserito (A21, B39, C81), insieme a C08 e B33, ma rendiamo il pacchetto più efficiente.', 'Lo sforzo totale resta ben entro il limite: 3474 ≤ 3683 (margine 209).', 'Se per te A82 è imprescindibile, alternativa: sostituire B39 con B31 e rimuovere C81; lo sforzo scende a 3412 (margine 271) mantenendo intatto il resto. Se preferisci questa variante, la formalizzo.'\}} \\
\texttt{RIFIUTO: \{'A11', 'B93', 'B09', 'C08', 'B33', 'A82', 'C85', 'C81', 'B39', 'A21'\}} \\
\texttt{PROPOSTA: \{'A11', 'B93', 'B09', 'C08', 'B33', 'C85', 'C81', 'B39', 'A21', 'B21'\}} \\
            }
        }
    }
    & \\ \\

    \theutterance \stepcounter{utterance}  
    & & & \multicolumn{4}{p{0.6\linewidth}}{
        \cellcolor[rgb]{0.9,0.9,0.9}{
            \makecell[{{p{\linewidth}}}]{
                \texttt{\tiny{[GM$\langle$P2]}}
                \texttt{ARGOMENTAZIONE: \{'Grazie: il tuo pacchetto è bilanciato e lascia buon margine. Per alzare ulteriormente l’efficienza mantenendo la stessa struttura portante, propongo uno scambio mirato in due mosse: (i) A11 + B21 → B96; (ii) B33 → A82.', 'Il risultato tiene tutti i tuoi pilastri appena inseriti (A21, B39, C81) insieme a C08, C85, B93 e B09, evita gli elementi che preferisci rimandare (nessun A60/B20) e resta entro il limite con margine di sicurezza: 3627 ≤ 3683 (margine 56).', 'Se per te B33 fosse imprescindibile, fallback immediato: teniamo B33 e non aggiungiamo A82, ma confermiamo lo scambio A11 + B21 → B96. Lo sforzo resta 3483 ≤ 3683 (margine 200). Fammi sapere quale preferisci e chiudiamo rapidamente.'\}} \\
\texttt{RIFIUTO: \{'A11', 'B93', 'B09', 'C08', 'B33', 'C85', 'C81', 'B39', 'A21', 'B21'\}} \\
\texttt{PROPOSTA: \{'A21', 'B93', 'B09', 'C08', 'C85', 'C81', 'B39', 'B96', 'A82'\}} \\
            }
        }
    }
     \\ \\

    \theutterance \stepcounter{utterance}  
    & & \multicolumn{4}{p{0.6\linewidth}}{
        \cellcolor[rgb]{0.9,0.9,0.9}{
            \makecell[{{p{\linewidth}}}]{
                \texttt{\tiny{[P1$\langle$GM]}}
                \texttt{ARGOMENTAZIONE: \{'Grazie: il tuo pacchetto è bilanciato e lascia buon margine. Per alzare ulteriormente l’efficienza mantenendo la stessa struttura portante, propongo uno scambio mirato in due mosse: (i) A11 + B21 → B96; (ii) B33 → A82.', 'Il risultato tiene tutti i tuoi pilastri appena inseriti (A21, B39, C81) insieme a C08, C85, B93 e B09, evita gli elementi che preferisci rimandare (nessun A60/B20) e resta entro il limite con margine di sicurezza: 3627 ≤ 3683 (margine 56).', 'Se per te B33 fosse imprescindibile, fallback immediato: teniamo B33 e non aggiungiamo A82, ma confermiamo lo scambio A11 + B21 → B96. Lo sforzo resta 3483 ≤ 3683 (margine 200). Fammi sapere quale preferisci e chiudiamo rapidamente.'\}} \\
\texttt{RIFIUTO: \{'A11', 'B93', 'B09', 'C08', 'B33', 'C85', 'C81', 'B39', 'A21', 'B21'\}} \\
\texttt{PROPOSTA: \{'A21', 'B93', 'B09', 'C08', 'C85', 'C81', 'B39', 'B96', 'A82'\}} \\
            }
        }
    }
    & \\ \\

    \theutterance \stepcounter{utterance}  
    & \multicolumn{4}{p{0.6\linewidth}}{
        \cellcolor[rgb]{0.9,0.9,0.9}{
            \makecell[{{p{\linewidth}}}]{
                \texttt{\tiny{[P1$\rangle$GM]}}
                \texttt{ARGOMENTAZIONE: \{'Mantengo intatta la tua struttura portante (B93, B09, C08, C81, B96, A82) e propongo tre scambi mirati per aumentare efficienza e margine: C85 → A11 (−79 sforzo), B39 → B33 (−25), A21 → B21 (+90). Saldo: −14 di sforzo totale.', 'Il set risultante resta ben entro il limite: 3613 ≤ 3683 (margine 70), con copertura più solida a parità di perimetro.', 'Se desideri mantenere A21, posso adattare 1‑per‑1 sostituendo B21 con A21 (sforzo 3523). Se invece vuoi tenere B39, posso offrire un’alternativa rimuovendo B21 e reinserendo B39 (sforzo 3517). Fammi sapere quale preferisci e la formalizzo.'\}} \\
\texttt{RIFIUTO: \{'A21', 'B93', 'B09', 'C08', 'C85', 'C81', 'B39', 'B96', 'A82'\}} \\
\texttt{PROPOSTA: \{'B93', 'B09', 'C08', 'C81', 'B96', 'A82', 'A11', 'B33', 'B21'\}} \\
            }
        }
    }
    & & \\ \\

    \theutterance \stepcounter{utterance}  
    & & \multicolumn{4}{p{0.6\linewidth}}{
        \cellcolor[rgb]{0.9,0.9,0.9}{
            \makecell[{{p{\linewidth}}}]{
                \texttt{\tiny{[GM$\rangle$P2]}}
                \texttt{ARGOMENTAZIONE: \{'Mantengo intatta la tua struttura portante (B93, B09, C08, C81, B96, A82) e propongo tre scambi mirati per aumentare efficienza e margine: C85 → A11 (−79 sforzo), B39 → B33 (−25), A21 → B21 (+90). Saldo: −14 di sforzo totale.', 'Il set risultante resta ben entro il limite: 3613 ≤ 3683 (margine 70), con copertura più solida a parità di perimetro.', 'Se desideri mantenere A21, posso adattare 1‑per‑1 sostituendo B21 con A21 (sforzo 3523). Se invece vuoi tenere B39, posso offrire un’alternativa rimuovendo B21 e reinserendo B39 (sforzo 3517). Fammi sapere quale preferisci e la formalizzo.'\}} \\
\texttt{RIFIUTO: \{'A21', 'B93', 'B09', 'C08', 'C85', 'C81', 'B39', 'B96', 'A82'\}} \\
\texttt{PROPOSTA: \{'B93', 'B09', 'C08', 'C81', 'B96', 'A82', 'A11', 'B33', 'B21'\}} \\
            }
        }
    }
    & \\ \\

    \theutterance \stepcounter{utterance}  
    & & & \multicolumn{4}{p{0.6\linewidth}}{
        \cellcolor[rgb]{0.9,0.9,0.9}{
            \makecell[{{p{\linewidth}}}]{
                \texttt{\tiny{[GM$\langle$P2]}}
                \texttt{ARGOMENTAZIONE: \{'Grazie per la proposta: ne condivido l’ossatura (B93, B09, C08, C81, B96, A82), che terrei intatta. Per aumentare efficienza e margine con un intervento minimo, propongo uno scambio 1‑per‑1: B21 → C85. Così riduciamo lo sforzo di 499 mantenendo copertura e coerenza con la tua enfasi sul margine.', 'Il nuovo set resta ampiamente entro il limite: 3114 ≤ 3683 (margine 569). Nessun inserimento di elementi che volevi evitare (A60, B20).', 'Se ti interessa un ulteriore affinamento a sforzo quasi invariato, posso valutare in un secondo momento anche B33 → B39 (+25 di sforzo; margine resterebbe molto ampio). Ma per ora mi concentro sullo scambio minimo proposto.'\}} \\
\texttt{PROPOSTA: \{'B93', 'B09', 'C08', 'C81', 'B96', 'A82', 'A11', 'B33', 'C85'\}} \\
            }
        }
    }
     \\ \\

    \theutterance \stepcounter{utterance}  
    & & \multicolumn{4}{p{0.6\linewidth}}{
        \cellcolor[rgb]{0.9,0.9,0.9}{
            \makecell[{{p{\linewidth}}}]{
                \texttt{\tiny{[P1$\langle$GM]}}
                \texttt{ARGOMENTAZIONE: \{'Grazie per la proposta: ne condivido l’ossatura (B93, B09, C08, C81, B96, A82), che terrei intatta. Per aumentare efficienza e margine con un intervento minimo, propongo uno scambio 1‑per‑1: B21 → C85. Così riduciamo lo sforzo di 499 mantenendo copertura e coerenza con la tua enfasi sul margine.', 'Il nuovo set resta ampiamente entro il limite: 3114 ≤ 3683 (margine 569). Nessun inserimento di elementi che volevi evitare (A60, B20).', 'Se ti interessa un ulteriore affinamento a sforzo quasi invariato, posso valutare in un secondo momento anche B33 → B39 (+25 di sforzo; margine resterebbe molto ampio). Ma per ora mi concentro sullo scambio minimo proposto.'\}} \\
\texttt{PROPOSTA: \{'B93', 'B09', 'C08', 'C81', 'B96', 'A82', 'A11', 'B33', 'C85'\}} \\
            }
        }
    }
    & \\ \\

    \theutterance \stepcounter{utterance}  
    & \multicolumn{4}{p{0.6\linewidth}}{
        \cellcolor[rgb]{0.9,0.9,0.9}{
            \makecell[{{p{\linewidth}}}]{
                \texttt{\tiny{[P1$\rangle$GM]}}
                \texttt{ARGOMENTAZIONE: \{'Intervento minimo e coerente con la tua struttura: sostituire B96 con B21. Così miglioriamo efficienza e aumentiamo il margine; lo sforzo totale scende a 3045 ≤ 3683 (margine 638).', 'Nessun altro cambiamento rispetto alla tua proposta: manteniamo B93, B09, C08, C81, A82, A11, B33 e C85 intatti.', 'Se preferisci un affinamento ulteriore sempre con ampio margine, posso proporre anche A82 → B31: sforzo 3164 (margine 519). Dimmi tu se vuoi questa variante.'\}} \\
\texttt{PROPOSTA: \{'B93', 'B09', 'C08', 'C81', 'B21', 'A82', 'A11', 'B33', 'C85'\}} \\
            }
        }
    }
    & & \\ \\

    \theutterance \stepcounter{utterance}  
    & & \multicolumn{4}{p{0.6\linewidth}}{
        \cellcolor[rgb]{0.9,0.9,0.9}{
            \makecell[{{p{\linewidth}}}]{
                \texttt{\tiny{[GM$\rangle$P2]}}
                \texttt{ARGOMENTAZIONE: \{'Intervento minimo e coerente con la tua struttura: sostituire B96 con B21. Così miglioriamo efficienza e aumentiamo il margine; lo sforzo totale scende a 3045 ≤ 3683 (margine 638).', 'Nessun altro cambiamento rispetto alla tua proposta: manteniamo B93, B09, C08, C81, A82, A11, B33 e C85 intatti.', 'Se preferisci un affinamento ulteriore sempre con ampio margine, posso proporre anche A82 → B31: sforzo 3164 (margine 519). Dimmi tu se vuoi questa variante.'\}} \\
\texttt{PROPOSTA: \{'B93', 'B09', 'C08', 'C81', 'B21', 'A82', 'A11', 'B33', 'C85'\}} \\
            }
        }
    }
    & \\ \\

    \theutterance \stepcounter{utterance}  
    & & & \multicolumn{4}{p{0.6\linewidth}}{
        \cellcolor[rgb]{0.9,0.9,0.9}{
            \makecell[{{p{\linewidth}}}]{
                \texttt{\tiny{[GM$\langle$P2]}}
                \texttt{ARGOMENTAZIONE: \{'Accolgo la tua struttura e propongo un unico scambio 1‑per‑1 che migliora l’efficacia senza intaccare il margine: B21 → B96.', 'Lo sforzo cresce solo di 69, restando ampiamente entro il limite: 3114 ≤ 3683 (margine 569). Nessun altro cambiamento: manteniamo B93, B09, C08, C81, A82, A11, B33 e C85 intatti, e continuiamo a evitare A60/B20.', 'Se preferisci comunque tenere B21, dimmelo: posso valutare un’alternativa minima che preservi il tuo perimetro. Ma questa sostituzione singola ci dà un pacchetto più solido con un margine ancora molto confortevole.'\}} \\
\texttt{RIFIUTO: \{'B93', 'B09', 'C08', 'C81', 'B21', 'A82', 'A11', 'B33', 'C85'\}} \\
\texttt{PROPOSTA: \{'B93', 'B09', 'C08', 'C81', 'B96', 'A82', 'A11', 'B33', 'C85'\}} \\
            }
        }
    }
     \\ \\

    \theutterance \stepcounter{utterance}  
    & & \multicolumn{4}{p{0.6\linewidth}}{
        \cellcolor[rgb]{0.9,0.9,0.9}{
            \makecell[{{p{\linewidth}}}]{
                \texttt{\tiny{[P1$\langle$GM]}}
                \texttt{ARGOMENTAZIONE: \{'Accolgo la tua struttura e propongo un unico scambio 1‑per‑1 che migliora l’efficacia senza intaccare il margine: B21 → B96.', 'Lo sforzo cresce solo di 69, restando ampiamente entro il limite: 3114 ≤ 3683 (margine 569). Nessun altro cambiamento: manteniamo B93, B09, C08, C81, A82, A11, B33 e C85 intatti, e continuiamo a evitare A60/B20.', 'Se preferisci comunque tenere B21, dimmelo: posso valutare un’alternativa minima che preservi il tuo perimetro. Ma questa sostituzione singola ci dà un pacchetto più solido con un margine ancora molto confortevole.'\}} \\
\texttt{RIFIUTO: \{'B93', 'B09', 'C08', 'C81', 'B21', 'A82', 'A11', 'B33', 'C85'\}} \\
\texttt{PROPOSTA: \{'B93', 'B09', 'C08', 'C81', 'B96', 'A82', 'A11', 'B33', 'C85'\}} \\
            }
        }
    }
    & \\ \\

    \theutterance \stepcounter{utterance}  
    & \multicolumn{4}{p{0.6\linewidth}}{
        \cellcolor[rgb]{0.9,0.9,0.9}{
            \makecell[{{p{\linewidth}}}]{
                \texttt{\tiny{[P1$\rangle$GM]}}
                \texttt{ARGOMENTAZIONE: \{'Accolgo la tua struttura e propongo un cambio minimo e mirato: sostituire C85 con B21. Manteniamo intatti tutti gli altri elementi, inclusi B96 e A82, e otteniamo un pacchetto più efficace con margine ancora sicuro.', 'Lo sforzo totale resta entro il limite: 3613 ≤ 3683 (margine 70).', 'Se preferisci preservare maggior margine, alternativa snella: lasciare C85 e sostituire B96 con B31; sforzo 3187 (margine 496). Dimmi quale strada preferisci e chiudiamo.'\}} \\
\texttt{RIFIUTO: \{'B93', 'B09', 'C08', 'C81', 'B96', 'A82', 'A11', 'B33', 'C85'\}} \\
\texttt{PROPOSTA: \{'B93', 'B09', 'C08', 'C81', 'B96', 'A82', 'A11', 'B33', 'B21'\}} \\
            }
        }
    }
    & & \\ \\

    \theutterance \stepcounter{utterance}  
    & & \multicolumn{4}{p{0.6\linewidth}}{
        \cellcolor[rgb]{0.9,0.9,0.9}{
            \makecell[{{p{\linewidth}}}]{
                \texttt{\tiny{[GM$\rangle$P2]}}
                \texttt{ARGOMENTAZIONE: \{'Accolgo la tua struttura e propongo un cambio minimo e mirato: sostituire C85 con B21. Manteniamo intatti tutti gli altri elementi, inclusi B96 e A82, e otteniamo un pacchetto più efficace con margine ancora sicuro.', 'Lo sforzo totale resta entro il limite: 3613 ≤ 3683 (margine 70).', 'Se preferisci preservare maggior margine, alternativa snella: lasciare C85 e sostituire B96 con B31; sforzo 3187 (margine 496). Dimmi quale strada preferisci e chiudiamo.'\}} \\
\texttt{RIFIUTO: \{'B93', 'B09', 'C08', 'C81', 'B96', 'A82', 'A11', 'B33', 'C85'\}} \\
\texttt{PROPOSTA: \{'B93', 'B09', 'C08', 'C81', 'B96', 'A82', 'A11', 'B33', 'B21'\}} \\
            }
        }
    }
    & \\ \\

    \theutterance \stepcounter{utterance}  
    & & & \multicolumn{4}{p{0.6\linewidth}}{
        \cellcolor[rgb]{0.9,0.9,0.9}{
            \makecell[{{p{\linewidth}}}]{
                \texttt{\tiny{[GM$\langle$P2]}}
                \texttt{ARGOMENTAZIONE: \{'Per chiudere rapidamente ripristino il set che avevi già formalizzato poco fa: stessa ossatura, ma senza gli elementi meno incisivi appena reintrodotti (A11, B33, B21).', 'Resta pienamente entro il limite: 3627 ≤ 3683 (margine 56). È una versione più efficace del tuo ultimo pacchetto a sforzo quasi identico.', 'Se per te B33 fosse imprescindibile, posso accettare B33 al posto di B39: sforzo 3602 (margine 81). Se invece preferisci mantenere B21, possiamo scambiare A21 → B21 e C85 → A11 per restare entro il limite: sforzo 3638 (margine 45). Dimmi quale preferisci.'\}} \\
\texttt{RIFIUTO: \{'B93', 'B09', 'C08', 'C81', 'B96', 'A82', 'A11', 'B33', 'B21'\}} \\
\texttt{PROPOSTA: \{'B93', 'B09', 'C08', 'C81', 'B96', 'A82', 'B39', 'A21', 'C85'\}} \\
            }
        }
    }
     \\ \\

    \theutterance \stepcounter{utterance}  
    & & \multicolumn{4}{p{0.6\linewidth}}{
        \cellcolor[rgb]{0.9,0.9,0.9}{
            \makecell[{{p{\linewidth}}}]{
                \texttt{\tiny{[P1$\langle$GM]}}
                \texttt{ARGOMENTAZIONE: \{'Per chiudere rapidamente ripristino il set che avevi già formalizzato poco fa: stessa ossatura, ma senza gli elementi meno incisivi appena reintrodotti (A11, B33, B21).', 'Resta pienamente entro il limite: 3627 ≤ 3683 (margine 56). È una versione più efficace del tuo ultimo pacchetto a sforzo quasi identico.', 'Se per te B33 fosse imprescindibile, posso accettare B33 al posto di B39: sforzo 3602 (margine 81). Se invece preferisci mantenere B21, possiamo scambiare A21 → B21 e C85 → A11 per restare entro il limite: sforzo 3638 (margine 45). Dimmi quale preferisci.'\}} \\
\texttt{RIFIUTO: \{'B93', 'B09', 'C08', 'C81', 'B96', 'A82', 'A11', 'B33', 'B21'\}} \\
\texttt{PROPOSTA: \{'B93', 'B09', 'C08', 'C81', 'B96', 'A82', 'B39', 'A21', 'C85'\}} \\
            }
        }
    }
    & \\ \\

    \theutterance \stepcounter{utterance}  
    & \multicolumn{4}{p{0.6\linewidth}}{
        \cellcolor[rgb]{0.9,0.9,0.9}{
            \makecell[{{p{\linewidth}}}]{
                \texttt{\tiny{[P1$\rangle$GM]}}
                \texttt{ARGOMENTAZIONE: \{'Rifiuto il tuo ultimo pacchetto perché mantiene elementi meno efficienti (B39, A21, B96) a scapito di soluzioni che avevamo già esplorato con buon esito.', 'Ti propongo di tornare alla spina dorsale condivisa (B93, B09, C08, A82) e ripristinare gli inserimenti che avevi già accettato (A11, B33, B21), sostituendo C81 con B31 per rafforzare la struttura. Restiamo ampiamente entro il limite: 3502 ≤ 3683 (margine 181).', 'Se per te serve più margine, posso adattare immediatamente B31 → C81 (sforzo 3164; margine 519) mantenendo inalterato il resto. In caso tu voglia tenere uno tra B96 o A21, possiamo ribilanciare 1‑per‑1 (es. B21 → B96) restando sotto soglia.', 'Pronto a chiudere su questa base; dimmi se preferisci la variante “alto margine” e la registro formalmente.'\}} \\
\texttt{RIFIUTO: \{'B93', 'B09', 'C08', 'C81', 'B96', 'A82', 'B39', 'A21', 'C85'\}} \\
\texttt{PROPOSTA: \{'A11', 'B33', 'C08', 'B93', 'B09', 'B21', 'B31', 'A82', 'C85'\}} \\
            }
        }
    }
    & & \\ \\

    \theutterance \stepcounter{utterance}  
    & & \multicolumn{4}{p{0.6\linewidth}}{
        \cellcolor[rgb]{0.9,0.9,0.9}{
            \makecell[{{p{\linewidth}}}]{
                \texttt{\tiny{[GM$\rangle$P2]}}
                \texttt{ARGOMENTAZIONE: \{'Rifiuto il tuo ultimo pacchetto perché mantiene elementi meno efficienti (B39, A21, B96) a scapito di soluzioni che avevamo già esplorato con buon esito.', 'Ti propongo di tornare alla spina dorsale condivisa (B93, B09, C08, A82) e ripristinare gli inserimenti che avevi già accettato (A11, B33, B21), sostituendo C81 con B31 per rafforzare la struttura. Restiamo ampiamente entro il limite: 3502 ≤ 3683 (margine 181).', 'Se per te serve più margine, posso adattare immediatamente B31 → C81 (sforzo 3164; margine 519) mantenendo inalterato il resto. In caso tu voglia tenere uno tra B96 o A21, possiamo ribilanciare 1‑per‑1 (es. B21 → B96) restando sotto soglia.', 'Pronto a chiudere su questa base; dimmi se preferisci la variante “alto margine” e la registro formalmente.'\}} \\
\texttt{RIFIUTO: \{'B93', 'B09', 'C08', 'C81', 'B96', 'A82', 'B39', 'A21', 'C85'\}} \\
\texttt{PROPOSTA: \{'A11', 'B33', 'C08', 'B93', 'B09', 'B21', 'B31', 'A82', 'C85'\}} \\
            }
        }
    }
    & \\ \\

    \theutterance \stepcounter{utterance}  
    & & & \multicolumn{4}{p{0.6\linewidth}}{
        \cellcolor[rgb]{0.9,0.9,0.9}{
            \makecell[{{p{\linewidth}}}]{
                \texttt{\tiny{[GM$\langle$P2]}}
                \texttt{ARGOMENTAZIONE: \{'Cambio minimo e coerente con le tue stesse varianti precedenti: A82 → C81. Manteniamo intatti B93, B09, C08, B21, B31, A11, B33 e C85.', 'Il set risulta più leggero e con ampio margine: 3164 ≤ 3683 (margine 519), senza introdurre elementi che avevi giudicato meno efficienti.', 'Se per te A82 fosse imprescindibile, alternativa 1‑per‑1: C08 → C81. In quel caso restiamo a 3491 ≤ 3683 (margine 192).', 'Preferisco chiudere sulla prima opzione (A82 → C81) per massimo margine e stabilità.'\}} \\
\texttt{RIFIUTO: \{'A11', 'B33', 'C08', 'B93', 'B09', 'B21', 'B31', 'A82', 'C85'\}} \\
\texttt{PROPOSTA: \{'A11', 'B33', 'C08', 'B93', 'B09', 'B21', 'B31', 'C85', 'C81'\}} \\
            }
        }
    }
     \\ \\

    \theutterance \stepcounter{utterance}  
    & & \multicolumn{4}{p{0.6\linewidth}}{
        \cellcolor[rgb]{0.9,0.9,0.9}{
            \makecell[{{p{\linewidth}}}]{
                \texttt{\tiny{[P1$\langle$GM]}}
                \texttt{ARGOMENTAZIONE: \{'Cambio minimo e coerente con le tue stesse varianti precedenti: A82 → C81. Manteniamo intatti B93, B09, C08, B21, B31, A11, B33 e C85.', 'Il set risulta più leggero e con ampio margine: 3164 ≤ 3683 (margine 519), senza introdurre elementi che avevi giudicato meno efficienti.', 'Se per te A82 fosse imprescindibile, alternativa 1‑per‑1: C08 → C81. In quel caso restiamo a 3491 ≤ 3683 (margine 192).', 'Preferisco chiudere sulla prima opzione (A82 → C81) per massimo margine e stabilità.'\}} \\
\texttt{RIFIUTO: \{'A11', 'B33', 'C08', 'B93', 'B09', 'B21', 'B31', 'A82', 'C85'\}} \\
\texttt{PROPOSTA: \{'A11', 'B33', 'C08', 'B93', 'B09', 'B21', 'B31', 'C85', 'C81'\}} \\
            }
        }
    }
    & \\ \\

    \theutterance \stepcounter{utterance}  
    & \multicolumn{4}{p{0.6\linewidth}}{
        \cellcolor[rgb]{0.9,0.9,0.9}{
            \makecell[{{p{\linewidth}}}]{
                \texttt{\tiny{[P1$\rangle$GM]}}
                \texttt{ARGOMENTAZIONE: \{'Il tuo set va bene e resta molto sotto soglia, ma con uno scambio minimo C81 → A82 rafforziamo il pacchetto senza rischi: lo sforzo totale passa a 3502 ≤ 3683, con margine 181.', 'Manteniamo intatti tutti gli altri elementi che hai selezionato (A11, B93, B09, C08, B33, B21, B31, C85) ed evitiamo comunque gli item più onerosi.', 'È una variante già emersa nella trattativa e coerente con la struttura condivisa; chiudiamo su questa versione snella ma più incisiva.'\}} \\
\texttt{RIFIUTO: \{'A11', 'B33', 'C08', 'B93', 'B09', 'B21', 'B31', 'C85', 'C81'\}} \\
\texttt{PROPOSTA: \{'A11', 'B33', 'C08', 'B93', 'B09', 'B21', 'B31', 'A82', 'C85'\}} \\
            }
        }
    }
    & & \\ \\

    \theutterance \stepcounter{utterance}  
    & & \multicolumn{4}{p{0.6\linewidth}}{
        \cellcolor[rgb]{0.9,0.9,0.9}{
            \makecell[{{p{\linewidth}}}]{
                \texttt{\tiny{[GM$\rangle$P2]}}
                \texttt{ARGOMENTAZIONE: \{'Il tuo set va bene e resta molto sotto soglia, ma con uno scambio minimo C81 → A82 rafforziamo il pacchetto senza rischi: lo sforzo totale passa a 3502 ≤ 3683, con margine 181.', 'Manteniamo intatti tutti gli altri elementi che hai selezionato (A11, B93, B09, C08, B33, B21, B31, C85) ed evitiamo comunque gli item più onerosi.', 'È una variante già emersa nella trattativa e coerente con la struttura condivisa; chiudiamo su questa versione snella ma più incisiva.'\}} \\
\texttt{RIFIUTO: \{'A11', 'B33', 'C08', 'B93', 'B09', 'B21', 'B31', 'C85', 'C81'\}} \\
\texttt{PROPOSTA: \{'A11', 'B33', 'C08', 'B93', 'B09', 'B21', 'B31', 'A82', 'C85'\}} \\
            }
        }
    }
    & \\ \\

    \theutterance \stepcounter{utterance}  
    & & & \multicolumn{4}{p{0.6\linewidth}}{
        \cellcolor[rgb]{0.9,0.9,0.9}{
            \makecell[{{p{\linewidth}}}]{
                \texttt{\tiny{[GM$\langle$P2]}}
                \texttt{ARGOMENTAZIONE: \{'Accolgo la tua struttura e propongo un unico scambio 1‑per‑1 a sforzo comparabile che migliora efficienza e margine: B21 → A21. Riduciamo lo sforzo di 90 mantenendo invariato tutto il resto: 3412 ≤ 3683 (margine 271).', 'È un intervento minimo, coerente con gli scambi mirati che hai proposto fin dall’inizio, e non tocca i tuoi pilastri (B93, B09, C08, B31, A82, A11, B33, C85).', 'Se per te A21 non fosse accettabile, segnalamelo: posso riconsiderare la tua versione precedente senza ulteriori modifiche.'\}} \\
\texttt{PROPOSTA: \{'A11', 'B33', 'C08', 'B93', 'B09', 'A21', 'B31', 'A82', 'C85'\}} \\
            }
        }
    }
     \\ \\

    \theutterance \stepcounter{utterance}  
    & & \multicolumn{4}{p{0.6\linewidth}}{
        \cellcolor[rgb]{0.9,0.9,0.9}{
            \makecell[{{p{\linewidth}}}]{
                \texttt{\tiny{[P1$\langle$GM]}}
                \texttt{ARGOMENTAZIONE: \{'Accolgo la tua struttura e propongo un unico scambio 1‑per‑1 a sforzo comparabile che migliora efficienza e margine: B21 → A21. Riduciamo lo sforzo di 90 mantenendo invariato tutto il resto: 3412 ≤ 3683 (margine 271).', 'È un intervento minimo, coerente con gli scambi mirati che hai proposto fin dall’inizio, e non tocca i tuoi pilastri (B93, B09, C08, B31, A82, A11, B33, C85).', 'Se per te A21 non fosse accettabile, segnalamelo: posso riconsiderare la tua versione precedente senza ulteriori modifiche.'\}} \\
\texttt{PROPOSTA: \{'A11', 'B33', 'C08', 'B93', 'B09', 'A21', 'B31', 'A82', 'C85'\}} \\
            }
        }
    }
    & \\ \\

    \theutterance \stepcounter{utterance}  
    & \multicolumn{4}{p{0.6\linewidth}}{
        \cellcolor[rgb]{0.9,0.9,0.9}{
            \makecell[{{p{\linewidth}}}]{
                \texttt{\tiny{[P1$\rangle$GM]}}
                \texttt{ARGOMENTAZIONE: \{'Apprezzo l’intervento minimo, ma A21 per me non è sostituibile a B21: ho bisogno esplicito di B21 per coprire una priorità chiave.', 'Propongo quindi il ripristino 1‑per‑1 A21 → B21: resta tutto identico e rimaniamo comodamente sotto soglia (3502 ≤ 3683, margine 181).', 'È l’aggiustamento più piccolo possibile per chiudere ora senza toccare gli altri pilastri che condividiamo (A11, B33, C08, B93, B09, B31, A82, C85).', 'Se per te servisse ulteriore margine, in seconda battuta posso valutare uno scambio alternativo su A82 mantenendo però B21 intatto.'\}} \\
\texttt{RIFIUTO: \{'A11', 'B33', 'C08', 'B93', 'B09', 'A21', 'B31', 'A82', 'C85'\}} \\
\texttt{PROPOSTA: \{'A11', 'B33', 'C08', 'B93', 'B09', 'B21', 'B31', 'A82', 'C85'\}} \\
            }
        }
    }
    & & \\ \\

    \theutterance \stepcounter{utterance}  
    & & \multicolumn{4}{p{0.6\linewidth}}{
        \cellcolor[rgb]{0.9,0.9,0.9}{
            \makecell[{{p{\linewidth}}}]{
                \texttt{\tiny{[GM$\rangle$P2]}}
                \texttt{ARGOMENTAZIONE: \{'Apprezzo l’intervento minimo, ma A21 per me non è sostituibile a B21: ho bisogno esplicito di B21 per coprire una priorità chiave.', 'Propongo quindi il ripristino 1‑per‑1 A21 → B21: resta tutto identico e rimaniamo comodamente sotto soglia (3502 ≤ 3683, margine 181).', 'È l’aggiustamento più piccolo possibile per chiudere ora senza toccare gli altri pilastri che condividiamo (A11, B33, C08, B93, B09, B31, A82, C85).', 'Se per te servisse ulteriore margine, in seconda battuta posso valutare uno scambio alternativo su A82 mantenendo però B21 intatto.'\}} \\
\texttt{RIFIUTO: \{'A11', 'B33', 'C08', 'B93', 'B09', 'A21', 'B31', 'A82', 'C85'\}} \\
\texttt{PROPOSTA: \{'A11', 'B33', 'C08', 'B93', 'B09', 'B21', 'B31', 'A82', 'C85'\}} \\
            }
        }
    }
    & \\ \\

    \theutterance \stepcounter{utterance}  
    & & & \multicolumn{4}{p{0.6\linewidth}}{
        \cellcolor[rgb]{0.9,0.9,0.9}{
            \makecell[{{p{\linewidth}}}]{
                \texttt{\tiny{[GM$\langle$P2]}}
                \texttt{ARGOMENTAZIONE: \{'Accolgo integralmente la tua struttura e applico l’unico scambio che avevi già proposto come variante “alto margine”: B31 → C81.', 'Teniamo fermi tutti i tuoi pilastri inclusi (A11, B33, C08, B93, B09, B21, A82, C85) e soddisfiamo la tua esigenza esplicita di mantenere B21.', 'Restiamo ampiamente sotto soglia: 3045 ≤ 3683 (margine 638). È l’aggiustamento più piccolo possibile per chiudere ora, senza introdurre elementi che avevi giudicato meno efficienti.'\}} \\
\texttt{PROPOSTA: \{'A11', 'B33', 'C08', 'B93', 'B09', 'B21', 'C81', 'A82', 'C85'\}} \\
            }
        }
    }
     \\ \\

    \theutterance \stepcounter{utterance}  
    & & \multicolumn{4}{p{0.6\linewidth}}{
        \cellcolor[rgb]{0.9,0.9,0.9}{
            \makecell[{{p{\linewidth}}}]{
                \texttt{\tiny{[P1$\langle$GM]}}
                \texttt{ARGOMENTAZIONE: \{'Accolgo integralmente la tua struttura e applico l’unico scambio che avevi già proposto come variante “alto margine”: B31 → C81.', 'Teniamo fermi tutti i tuoi pilastri inclusi (A11, B33, C08, B93, B09, B21, A82, C85) e soddisfiamo la tua esigenza esplicita di mantenere B21.', 'Restiamo ampiamente sotto soglia: 3045 ≤ 3683 (margine 638). È l’aggiustamento più piccolo possibile per chiudere ora, senza introdurre elementi che avevi giudicato meno efficienti.'\}} \\
\texttt{PROPOSTA: \{'A11', 'B33', 'C08', 'B93', 'B09', 'B21', 'C81', 'A82', 'C85'\}} \\
            }
        }
    }
    & \\ \\

    \theutterance \stepcounter{utterance}  
    & \multicolumn{4}{p{0.6\linewidth}}{
        \cellcolor[rgb]{0.9,0.9,0.9}{
            \makecell[{{p{\linewidth}}}]{
                \texttt{\tiny{[P1$\rangle$GM]}}
                \texttt{ARGOMENTAZIONE: \{'Il tuo set è solido e con ampio margine; propongo un aggiustamento minimo e 1‑per‑1 per rafforzarlo ulteriormente: C81 → B31.', 'Restiamo comodamente sotto il limite: 3502 ≤ 3683 (margine 181). Manteniamo intatti tutti gli altri elementi che hai scelto (A11, B33, C08, B93, B09, B21, A82, C85) e continuiamo a evitare gli item più onerosi.', 'È la variante già emersa nella nostra trattativa quando la struttura era invariata: stesso perimetro, ma con maggiore solidità operativa.'\}} \\
\texttt{PROPOSTA: \{'A11', 'B33', 'C08', 'B93', 'B09', 'B21', 'B31', 'A82', 'C85'\}} \\
            }
        }
    }
    & & \\ \\

    \theutterance \stepcounter{utterance}  
    & & \multicolumn{4}{p{0.6\linewidth}}{
        \cellcolor[rgb]{0.9,0.9,0.9}{
            \makecell[{{p{\linewidth}}}]{
                \texttt{\tiny{[GM$\rangle$P2]}}
                \texttt{ARGOMENTAZIONE: \{'Il tuo set è solido e con ampio margine; propongo un aggiustamento minimo e 1‑per‑1 per rafforzarlo ulteriormente: C81 → B31.', 'Restiamo comodamente sotto il limite: 3502 ≤ 3683 (margine 181). Manteniamo intatti tutti gli altri elementi che hai scelto (A11, B33, C08, B93, B09, B21, A82, C85) e continuiamo a evitare gli item più onerosi.', 'È la variante già emersa nella nostra trattativa quando la struttura era invariata: stesso perimetro, ma con maggiore solidità operativa.'\}} \\
\texttt{PROPOSTA: \{'A11', 'B33', 'C08', 'B93', 'B09', 'B21', 'B31', 'A82', 'C85'\}} \\
            }
        }
    }
    & \\ \\

    \theutterance \stepcounter{utterance}  
    & & & \multicolumn{4}{p{0.6\linewidth}}{
        \cellcolor[rgb]{0.9,0.9,0.9}{
            \makecell[{{p{\linewidth}}}]{
                \texttt{\tiny{[GM$\langle$P2]}}
                \texttt{ARGOMENTAZIONE: \{'Il tuo set è coerente; per chiudere propongo di fissare la variante alto‑margine già sul tavolo: scambio 1‑per‑1 B31 → C81, tutto il resto invariato.', 'Così aumentiamo nettamente il margine senza cambiare perimetro né introdurre elementi che volevi evitare: 3045 ≤ 3683 (margine 638). B21 resta garantito.', 'È l’aggiustamento minimo per massimizzare stabilità operativa. Se ti sta bene, chiudiamo subito; se preferisci un micro‑adattamento diverso, indicami l’unico scambio che vuoi e lo registro.'\}} \\
\texttt{RIFIUTO: \{'A11', 'B33', 'C08', 'B93', 'B09', 'B21', 'B31', 'A82', 'C85'\}} \\
\texttt{PROPOSTA: \{'A11', 'B33', 'C08', 'B93', 'B09', 'B21', 'C81', 'A82', 'C85'\}} \\
            }
        }
    }
     \\ \\

    \theutterance \stepcounter{utterance}  
    & & \multicolumn{4}{p{0.6\linewidth}}{
        \cellcolor[rgb]{0.9,0.9,0.9}{
            \makecell[{{p{\linewidth}}}]{
                \texttt{\tiny{[P1$\langle$GM]}}
                \texttt{ARGOMENTAZIONE: \{'Il tuo set è coerente; per chiudere propongo di fissare la variante alto‑margine già sul tavolo: scambio 1‑per‑1 B31 → C81, tutto il resto invariato.', 'Così aumentiamo nettamente il margine senza cambiare perimetro né introdurre elementi che volevi evitare: 3045 ≤ 3683 (margine 638). B21 resta garantito.', 'È l’aggiustamento minimo per massimizzare stabilità operativa. Se ti sta bene, chiudiamo subito; se preferisci un micro‑adattamento diverso, indicami l’unico scambio che vuoi e lo registro.'\}} \\
\texttt{RIFIUTO: \{'A11', 'B33', 'C08', 'B93', 'B09', 'B21', 'B31', 'A82', 'C85'\}} \\
\texttt{PROPOSTA: \{'A11', 'B33', 'C08', 'B93', 'B09', 'B21', 'C81', 'A82', 'C85'\}} \\
            }
        }
    }
    & \\ \\

    \theutterance \stepcounter{utterance}  
    & \multicolumn{4}{p{0.6\linewidth}}{
        \cellcolor[rgb]{0.9,0.9,0.9}{
            \makecell[{{p{\linewidth}}}]{
                \texttt{\tiny{[P1$\rangle$GM]}}
                \texttt{ARGOMENTAZIONE: \{'Intervento minimo 1‑per‑1: C81 → B31. Rafforza il pacchetto mantenendo intatta l’ossatura e restando comodamente sotto soglia (3502 ≤ 3683; margine 181).', 'Nessun altro cambiamento: teniamo fermi A11, B33, C08, B93, B09, B21, A82 e C85; continuiamo a evitare gli elementi più onerosi già esclusi.', 'Se per te il margine extra è imprescindibile, posso confermare anche la tua versione con C81; in assenza di contrarie indicazioni, preferisco chiudere con B31 per maggiore solidità.'\}} \\
\texttt{RIFIUTO: \{'A11', 'B33', 'C08', 'B93', 'B09', 'B21', 'C81', 'A82', 'C85'\}} \\
\texttt{PROPOSTA: \{'A11', 'B33', 'C08', 'B93', 'B09', 'B21', 'B31', 'A82', 'C85'\}} \\
            }
        }
    }
    & & \\ \\

    \theutterance \stepcounter{utterance}  
    & & \multicolumn{4}{p{0.6\linewidth}}{
        \cellcolor[rgb]{0.9,0.9,0.9}{
            \makecell[{{p{\linewidth}}}]{
                \texttt{\tiny{[GM$\rangle$P2]}}
                \texttt{ARGOMENTAZIONE: \{'Intervento minimo 1‑per‑1: C81 → B31. Rafforza il pacchetto mantenendo intatta l’ossatura e restando comodamente sotto soglia (3502 ≤ 3683; margine 181).', 'Nessun altro cambiamento: teniamo fermi A11, B33, C08, B93, B09, B21, A82 e C85; continuiamo a evitare gli elementi più onerosi già esclusi.', 'Se per te il margine extra è imprescindibile, posso confermare anche la tua versione con C81; in assenza di contrarie indicazioni, preferisco chiudere con B31 per maggiore solidità.'\}} \\
\texttt{RIFIUTO: \{'A11', 'B33', 'C08', 'B93', 'B09', 'B21', 'C81', 'A82', 'C85'\}} \\
\texttt{PROPOSTA: \{'A11', 'B33', 'C08', 'B93', 'B09', 'B21', 'B31', 'A82', 'C85'\}} \\
            }
        }
    }
    & \\ \\

    \theutterance \stepcounter{utterance}  
    & & & \multicolumn{4}{p{0.6\linewidth}}{
        \cellcolor[rgb]{0.9,0.9,0.9}{
            \makecell[{{p{\linewidth}}}]{
                \texttt{\tiny{[GM$\langle$P2]}}
                \texttt{ARGOMENTAZIONE: \{'Condivido la tua preferenza per chiudere con questa struttura: copre bene le priorità, resta sotto soglia con margine adeguato (3502 ≤ 3683) ed evita gli elementi più onerosi.', 'Manteniamo tutti i pilastri condivisi (A11, B33, C08, B93, B09, B21, B31, A82, C85) senza complicare ulteriormente lo scambio.'\}} \\
\texttt{ACCORDO: \{'A11', 'B33', 'C08', 'B93', 'B09', 'B21', 'B31', 'A82', 'C85'\}} \\
            }
        }
    }
     \\ \\

    \theutterance \stepcounter{utterance}  
    & & & \multicolumn{2}{p{0.3\linewidth}}{
        \cellcolor[rgb]{0.9,0.9,0.9}{
            \makecell[{{p{\linewidth}}}]{
                \texttt{\tiny{[GM$|$GM]}}
                \texttt{game successful} \\
            }
        }
    }
    & & \\ \\

    \theutterance \stepcounter{utterance}  
    & & & \multicolumn{2}{p{0.3\linewidth}}{
        \cellcolor[rgb]{0.9,0.9,0.9}{
            \makecell[{{p{\linewidth}}}]{
                \texttt{\tiny{[GM$|$GM]}}
                \texttt{end game} \\
            }
        }
    }
    & & \\ \\

\end{supertabular}
}

\end{document}
